\documentclass[11pt,a4paper]{article}
\usepackage{fontspec,amsmath,amssymb,amsthm,microtype}
\setmainfont{lmroman10-regular.otf}[BoldFont=lmroman10-bold.otf,ItalicFont=lmroman10-italic.otf,BoldItalicFont=lmroman10-bolditalic.otf]
\usepackage{xcolor}
\usepackage[most]{tcolorbox}
\usepackage{mathrsfs}
\providecommand{\lt}{<}
\providecommand{\gt}{>}
\setmonofont{DejaVuSansMono.ttf}[BoldFont=DejaVuSansMono-Bold.ttf,ItalicFont=DejaVuSansMono-Oblique.ttf,BoldItalicFont=DejaVuSansMono-BoldOblique.ttf,Scale=0.86]
\newfontfamily\leanfallback{STIXTwoMath-Regular.otf}[Scale=0.86]
\definecolor{leanrule}{gray}{0.62}
\definecolor{leanback}{gray}{0.962}
\definecolor{leantitle}{gray}{0.22}
\definecolor{leanmuted}{gray}{0.38}
\newtcolorbox{leanbox}[1]{enhanced,breakable,colback=leanback,colframe=leanrule,
  boxrule=0.4pt,arc=0.6mm,left=2.4mm,right=2.4mm,top=1.8mm,bottom=1.8mm,
  before skip=0.8em,after skip=1em,fonttitle=\small\bfseries,coltitle=leantitle,
  colbacktitle=leanback,titlerule=0.3pt,title={#1}}
\newlength{\leancharwidth}
\newenvironment{leancode}{\par\smallskip\begingroup\ttfamily\footnotesize\raggedright
  \settowidth{\leancharwidth}{M}\setlength{\parskip}{0pt}\setlength{\parindent}{0pt}
  \frenchspacing}{\par\endgroup\smallskip}
\newcommand{\leanline}[2]{\hangindent=\dimexpr#1\leancharwidth+4\leancharwidth\relax
  \hangafter=1\noindent\hspace*{#1\leancharwidth}#2\par}
\newcommand{\leansource}[1]{\par{\footnotesize\color{leanmuted}#1}\par}
\newcommand{\leanlabel}[1]{\par\smallskip\noindent{\small\bfseries #1}\par}
\usepackage[margin=25mm]{geometry}
\usepackage[unicode,hidelinks]{hyperref}
\hypersetup{pdfauthor={LangevinLean Project}}
\setlength{\emergencystretch}{3em}
\setlength{\parskip}{0.4em}
\setlength{\parindent}{0pt}
\allowdisplaybreaks[2]
\raggedbottom
\title{\texorpdfstring{The Gaussian marginal and primitive of the free Ornstein–Uhlenbeck velocity}{The Gaussian marginal and primitive of the free Ornstein–Uhlenbeck velocity}}
\author{LangevinLean Project}
\date{}
\begin{document}\maketitle
Draft mathematical exposition; these notes may contain mistakes.


The free Ornstein–Uhlenbeck velocity is the solution of the kinetic velocity equation without force: started at \(V\) and driven by \(\sqrt{2\gamma}\,dW\), it equals


\[U_t=\psi_0(t)V+N_t,\qquad\psi_0(t)=e^{-\gamma t},\qquad N_t=\sqrt{2\gamma}\int_0^te^{-\gamma(t-s)}\,dW_s.\]


It appears in the second-order algorithm, where the Hessian at the initial position multiplies \(\int_0^t\int_0^sU_r\,dr\,ds\), the part of the displacement that is linear in the initial velocity and in the fresh noise. This section treats \(U\) as a path in \(L^2(\mathbb P;E)\) for a given initial velocity \(V\) and a given continuous innovation path \(N\). We prove: its primitive is \(\int_0^tU_s\,ds=\psi_1(t)V+\int_0^tN_s\,ds\) with \(\psi_1(t)=\int_0^t\psi_0\); and if \(V\sim\mathcal N(0,I)\) and \(N_t\) is independent of \(V\) with law \(\mathcal N(0,(1-\psi_0(t)^2)I)\), then \(U_t\sim\mathcal N(0,I)\). The independence hypothesis is verified later for the Brownian reference process, from the independence of past and future Brownian increments.


\section{\texorpdfstring{An independent thermal innovation preserves the standard Gaussian law}{An independent thermal innovation preserves the standard Gaussian law}}


\textbf{Claim.} Let \(E\) be a finite-dimensional real inner product space with its Borel \(\sigma\)-algebra and standard Gaussian law \(\mathcal N=\mathcal N(0,I)\). Let \(a\in\mathbb R\) with \(a^2\le1\), and let \(\mathcal N_c\) be the image of \(\mathcal N\) under \(z\mapsto cz\), \(c=\sqrt{1-a^2}\). Then the image of \(\mathcal N\otimes\mathcal N_c\) under \((x,y)\mapsto ax+y\) is \(\mathcal N\). In terms of random vectors: if \(Z\sim\mathcal N\) and \(Y\sim\mathcal N_c\) are independent, then \(aZ+Y\sim\mathcal N\).




\textbf{Proof.} The product of two image measures is the image of the product measure under the product map: \(\mathcal N\otimes\mathcal N_c\) is the image of \(\mathcal N\otimes\mathcal N\) under \((z,z')\mapsto(z,cz')\). Composing the two maps, the measure in question is the image of \(\mathcal N\otimes\mathcal N\) under \((z,z')\mapsto az+cz'\). Since \(a^2+c^2=a^2+(1-a^2)=1\), the earlier linear-combination lemma for independent standard Gaussians applies: for independent \(Z,Z'\sim\mathcal N\) and \(a^2+c^2=1\), \(aZ+cZ'\sim\mathcal N\). (Its proof compares characteristic functions: \(\mathbb E\,e^{i\langle\xi,aZ+cZ'\rangle}=e^{-a^2\|\xi\|^2/2}e^{-c^2\|\xi\|^2/2}=e^{-\|\xi\|^2/2}\), and a law on \(E\) is determined by its characteristic function.) \(\square\)




The conclusion is an identity of probability measures, not a statement about covariances only.

\begin{leanbox}{Lean: \texttt{stdGaussian\_\allowbreak{}add\_\allowbreak{}thermal}}
\leansource{\texttt{stdGaussian\_\allowbreak{}add\_\allowbreak{}thermal} \quad KineticLangevin/\allowbreak{}FreeOU.lean, lines 15–24}
\begin{leancode}
\leanline{0}{theorem\ stdGaussian\_add\_thermal\ (a\ :\ ℝ)\ (ha\ :\ a\textasciicircum{}2≤1)\ :}
\leanline{4}{((stdGaussian\ E).prod\ ((stdGaussian\ E).map\ (fun\ z\ ↦\ Real.sqrt\ (1-a\textasciicircum{}2)\ •\ z))).map}
\leanline{6}{(fun\ p\ ↦\ a\ •\ p.1+p.2)\ =\ stdGaussian\ E}
\end{leancode}
\leanlabel{How the Lean proof follows the argument}
In \texttt{std\allowbreak{}Gaussian\_\allowbreak{}add\_\allowbreak{}thermal}, the law \(\mathcal N_c\) is written \texttt{(std\allowbreak{}Gaussian E).\allowbreak{}map (fun z \ensuremath{\mapsto} Real.\allowbreak{}sqrt (1-\allowbreak{}a\textasciicircum{}2) \ensuremath{\bullet} z)} and the hypothesis is \texttt{ha : a\textasciicircum{}2\ensuremath{\le}1}. \texttt{hp} is \texttt{Measure.\allowbreak{}map\_\allowbreak{}prod\_\allowbreak{}map}, the identity between a product of image measures and the image of the product measure; \texttt{Measure.\allowbreak{}map\_\allowbreak{}map} composes the two maps. The goal is then the earlier lemma \texttt{std\allowbreak{}Gaussian\_\allowbreak{}linear\_\allowbreak{}combination} with coefficients \(a\) and \(\sqrt{1-a^2}\); its hypothesis \(a^2+(\sqrt{1-a^2})^2=1\) is proved by \texttt{Real.\allowbreak{}sq\_\allowbreak{}sqrt} (using \(1-a^2\ge0\), from \texttt{sub\_\allowbreak{}nonneg.\allowbreak{}mpr ha}) and \texttt{ring}.
\end{leanbox}

\section{\texorpdfstring{The free Ornstein–Uhlenbeck path and its primitive}{The free Ornstein–Uhlenbeck path and its primitive}}


\textbf{Definition.} Let \((\Omega,\mathbb P)\) be a measure space, \(E\) a real normed space, \(\gamma\in\mathbb R\), \(V\in L^2(\mathbb P;E)\) and \(N:\mathbb R\to L^2(\mathbb P;E)\). The free Ornstein–Uhlenbeck velocity is the path


\[U_t=\psi_0(t)\,V+N_t\in L^2(\mathbb P;E),\qquad\psi_0(t)=e^{-\gamma t}.\]


\textbf{Claim.} (a) If \(N\) is continuous, then \(t\mapsto U_t\) is continuous. (b) If moreover \(E\) is complete, then for every \(t\in\mathbb R\),


\[\int_0^tU_s\,ds=\psi_1(t)\,V+\int_0^tN_s\,ds,\qquad\psi_1(t)=\int_0^t\psi_0(s)\,ds,\]


with Bochner integrals in \(L^2(\mathbb P;E)\) (oriented integrals when \(t<0\)).




\textbf{Proof.} (a) \(t\mapsto\psi_0(t)\) is continuous, scalar multiplication \(\mathbb R\times L^2\to L^2\) is continuous, and the sum of continuous paths is continuous.




(b) Both \(s\mapsto\psi_0(s)V\) and \(s\mapsto N_s\) are continuous, hence integrable on the interval between \(0\) and \(t\). By linearity of the integral, \(\int_0^tU_s\,ds=\int_0^t\psi_0(s)V\,ds+\int_0^tN_s\,ds\). For a fixed vector \(V\), \(\int_0^t\psi_0(s)V\,ds=\bigl(\int_0^t\psi_0(s)\,ds\bigr)V=\psi_1(t)V\), which is the definition of \(\psi_1\). \(\square\)




For the Brownian reference process, \(N_t=\sqrt{2\gamma}\int_0^te^{-\gamma(t-s)}dW_s\) and the second term is the position innovation \(\sqrt{2\gamma}\int_0^t\psi_1(t-s)\,dW_s\). The initial-velocity part and the noise part of the displacement are thus integrated on the same Brownian interval.

\begin{leanbox}{Lean: \texttt{freeOUVelocity}, \texttt{continuous\_\allowbreak{}freeOUVelocity}, \texttt{freeOUVelocity\_\allowbreak{}primitive}}
\leansource{\texttt{freeOUVelocity} \quad KineticLangevin/\allowbreak{}FreeOU.lean, lines 32–33}
\begin{leancode}
\leanline{0}{def\ freeOUVelocity\ (γ\ :\ ℝ)\ (V\ :\ Lp\ E\ 2\ P)\ (N\ :\ ℝ→Lp\ E\ 2\ P)\ (t\ :\ ℝ)\ :\ Lp\ E\ 2\ P\ :=}
\leanline{2}{psi\ γ\ 0\ t\ •\ V+N\ t}
\end{leancode}
\leansource{\texttt{continuous\_\allowbreak{}freeOUVelocity} \quad KineticLangevin/\allowbreak{}FreeOU.lean, lines 35–37}
\begin{leancode}
\leanline{0}{theorem\ continuous\_freeOUVelocity\ (γ\ :\ ℝ)\ (V\ :\ Lp\ E\ 2\ P)\ (N\ :\ ℝ→Lp\ E\ 2\ P)}
\leanline{4}{(hN\ :\ Continuous\ N)\ :\ Continuous\ (freeOUVelocity\ P\ γ\ V\ N)}
\end{leancode}
\leansource{\texttt{freeOUVelocity\_\allowbreak{}primitive} \quad KineticLangevin/\allowbreak{}FreeOU.lean, lines 40–47}
\begin{leancode}
\leanline{0}{theorem\ freeOUVelocity\_primitive\ [CompleteSpace\ E]\ (γ\ :\ ℝ)\ (V\ :\ Lp\ E\ 2\ P)\ (N\ :\ ℝ→Lp\ E\ 2\ P)}
\leanline{4}{(hN\ :\ Continuous\ N)\ (t\ :\ ℝ)\ :}
\leanline{4}{(∫\ s\ in\ 0..t,freeOUVelocity\ P\ γ\ V\ N\ s)=psi\ γ\ 1\ t\ •\ V+∫\ s\ in\ 0..t,N\ s}
\end{leancode}
\leanlabel{How the Lean proof follows the argument}
\texttt{free\allowbreak{}OUVelocity P \ensuremath{\gamma} V N t} is defined as \texttt{psi \ensuremath{\gamma} 0 t \ensuremath{\bullet} V+N t} in \texttt{Lp E 2 P}, where \texttt{psi \ensuremath{\gamma} 0 t} is \(e^{-\gamma t}\) and \texttt{psi \ensuremath{\gamma} 1} is its primitive, defined by the interval integral from \(0\). \texttt{continuous\_\allowbreak{}free\allowbreak{}OUVelocity} combines \texttt{continuous\_\allowbreak{}psi \ensuremath{\gamma} 0}, scalar multiplication by the constant \texttt{V}, and \texttt{h\allowbreak{}N : Continuous N}.




\texttt{free\allowbreak{}OUVelocity\_\allowbreak{}primitive} (with \texttt{[Complete\allowbreak{}Space E]}) splits the interval integral with \texttt{interval\allowbreak{}Integral.\allowbreak{}integral\_\allowbreak{}add}, using interval integrability \texttt{hi} of the first term and that of \texttt{N}; \texttt{interval\allowbreak{}Integral.\allowbreak{}integral\_\allowbreak{}smul\_\allowbreak{}const} moves the constant vector \texttt{V} outside, and the remaining scalar integral is \texttt{psi \ensuremath{\gamma} 1 t} by definition (\texttt{rfl}).
\end{leanbox}

\section{\texorpdfstring{The Gaussian marginal of the free velocity}{The Gaussian marginal of the free velocity}}


\textbf{Claim.} Let \(E\) be a finite-dimensional real inner product space, \((\Omega,\mathbb P)\) a measure space, \(\gamma\ge0\), \(t\ge0\), \(V\in L^2(\mathbb P;E)\) and \(N:\mathbb R\to L^2(\mathbb P;E)\). Suppose the pair \((V,N_t)\) has the law


\[\mathcal N(0,I)\otimes\bigl(\text{image of }\mathcal N(0,I)\text{ under }z\mapsto\sqrt{1-\psi_0(t)^2}\,z\bigr),\]


that is, \(V\sim\mathcal N(0,I)\), \(N_t\sim\mathcal N(0,(1-\psi_0(t)^2)I)\), and \(V\), \(N_t\) are independent. Then \(U_t=\psi_0(t)V+N_t\) has law \(\mathcal N(0,I)\).




\textbf{Proof.} Since \(\gamma\ge0\) and \(t\ge0\), \(0\le\psi_0(t)=e^{-\gamma t}\le1\), so \(\psi_0(t)^2\le1\). By the first step with \(a=\psi_0(t)\), the image of the pair law under \((x,y)\mapsto\psi_0(t)x+y\) is \(\mathcal N(0,I)\). The law of a measurable function of the pair \((V,N_t)\) is the image of the pair law under that function, so \(\psi_0(t)V+N_t\), computed pointwise, has law \(\mathcal N(0,I)\). Finally, the \(L^2\) class \(U_t\) is represented almost everywhere by \(\omega\mapsto\psi_0(t)V(\omega)+N_t(\omega)\), and almost everywhere equal random vectors have the same law. \(\square\)




The hypothesis involves only the joint law of \((V,N_t)\). In the second-order comparison, \(V\) is the stationary velocity at the start of a step and \(N_t\) is built from the Brownian increments after that time; the required independence follows from the independence of past and future Brownian increments, proved later. No independence between \(U_t\) and the position discrepancy of the coupled algorithm is needed.

\begin{leanbox}{Lean: \texttt{freeOUVelocity\_\allowbreak{}hasLaw}}
\leansource{\texttt{freeOUVelocity\_\allowbreak{}hasLaw} \quad KineticLangevin/\allowbreak{}FreeOU.lean, lines 51–70}
\begin{leancode}
\leanline{0}{theorem\ freeOUVelocity\_hasLaw\ \{Ω\ E\ :\ Type*\}\ [MeasurableSpace\ Ω]}
\leanline{4}{[NormedAddCommGroup\ E]\ [InnerProductSpace\ ℝ\ E]\ [FiniteDimensional\ ℝ\ E]}
\leanline{4}{[MeasurableSpace\ E]\ [BorelSpace\ E]}
\leanline{4}{(P\ :\ Measure\ Ω)\ (γ\ :\ ℝ)\ (hγ\ :\ 0≤γ)\ (V\ :\ Lp\ E\ 2\ P)\ (N\ :\ ℝ→Lp\ E\ 2\ P)}
\leanline{4}{(t\ :\ ℝ)\ (ht\ :\ 0≤t)}
\leanline{4}{(hpair\ :\ HasLaw\ (fun\ ω\ ↦\ (V\ ω,N\ t\ ω))}
\leanline{6}{((stdGaussian\ E).prod\ ((stdGaussian\ E).map\ (fun\ z\ ↦\ Real.sqrt\ (1-(psi\ γ\ 0\ t)\textasciicircum{}2)\ •\ z)))\ P)\ :}
\leanline{4}{HasLaw\ (freeOUVelocity\ P\ γ\ V\ N\ t)\ (stdGaussian\ E)\ P}
\end{leancode}
\leanlabel{How the Lean proof follows the argument}
In \texttt{free\allowbreak{}OUVelocity\_\allowbreak{}has\allowbreak{}Law}, the hypothesis \texttt{hpair} is a \texttt{Has\allowbreak{}Law} statement for the pair \texttt{fun \ensuremath{\omega} \ensuremath{\mapsto} (V \ensuremath{\omega},N t \ensuremath{\omega})} with the product law written as above, and the conclusion is \texttt{Has\allowbreak{}Law (free\allowbreak{}OUVelocity P \ensuremath{\gamma} V N t) (std\allowbreak{}Gaussian E) P}.



\begin{enumerate}


\item \texttt{hp} is \(\psi_0(t)^2\le1\), from \texttt{psi\_\allowbreak{}nonneg} and \texttt{psi\_\allowbreak{}zero\_\allowbreak{}le\_\allowbreak{}one} (which uses \texttt{h\ensuremath{\gamma}} and \texttt{ht}) by \texttt{nlinarith}.

\item \texttt{hsum} states that the map \((x,y)\mapsto\psi_0(t)x+y\) sends the pair law to \(\mathcal N(0,I)\); it is built from measurability (\texttt{fun\_\allowbreak{}prop}) and \texttt{std\allowbreak{}Gaussian\_\allowbreak{}add\_\allowbreak{}thermal}.

\item \texttt{hsum.\allowbreak{}comp hpair} composes the two laws, and \texttt{.\allowbreak{}congr} replaces the pointwise expression by the \(L^2\) class almost everywhere; the almost-everywhere identity comes from \texttt{Lp.\allowbreak{}coe\allowbreak{}Fn\_\allowbreak{}add} and \texttt{Lp.\allowbreak{}coe\allowbreak{}Fn\_\allowbreak{}smul}.

\end{enumerate}
\end{leanbox}

\end{document}
